\section{Technologies Used and Techniques Applied}
\label{sec:technologies-used-and-techniques-applied}

\subsection{Choice of Gait-Recognition Approach}
\label{subsec:choice-of-gait-recognition-approach}

Two main forms of \ac{gr} are discussed in Chapter \ref{chap:context} and \ref{chap:new-ideas}:
Appearance (e.g., silhouette, GEI), and Model-based (e.g., 2 or 3-dimensional skeletons).
To settle on the most suitable approach to use for this project (not just for efficacy;
the project needed to be completed within the given timeframe and able to operate on available hardware)
benchmarks were performed to evaluate the technologies required to implement either approach.

The same system was used to perform all benchmarking, ensuring results could be compared directly and would represent system performance.
As a target of this project is for the solution to be efficient enough to run on 'consumer'-grade (low-medium cost and power) hardware,
the system used for these tests --- an LG 16Z90P laptop --- was of no significant cost, and was only equipped with an integrated \ac{gpu}.

Specifications of this system relevant to performance are detailed in Table \ref{tab:benchmark-system-specs} below.

\begin{table}[H]
    \begin{scriptsize}
        \centering
        \caption{Benchmark and Testing System Specifications}
        \label{tab:benchmark-system-specs}
        \begin{tabular}{|L{0.3\textwidth}|L{0.6\textwidth}|}
            \hline
            \textbf{Component}        & \textbf{Specification}                                                           \\ \hline
            \textbf{Processor}        & 11th Gen Intel Core i7-1165G7 @ 2.80GHz, 2803 Mhz, 4 Cores, 8 Logical Processors \\ \hline
            \textbf{Physical Memory}  & 16.0 GB                                                                          \\ \hline
            \textbf{Graphics}         & Intel Iris Xe Graphics (CPU-integrated)                                          \\ \hline
            \textbf{Operating System} & Microsoft Windows 11 Home 64-bit                                                 \\ \hline
        \end{tabular}
    \end{scriptsize}
\end{table}

A single data sample was used for benchmarking --- a set of 106 lateral movement frames at 1080p resolution taken from the publically accessible GRIDDS sample dataset \parencite{nunesGRIDDSGaitRecognition2019} --- with both approaches applied to the same sample.
The aim here was to measure the time taken by each approach to perform feature extraction for an entire movement cycle (i.e., a person walking through the camera's field of view laterally).

\paragraph{Model-based \ac{gr}} The OpenPose framework was used here. It enables extraction of 2-dimensional (or 3-dimensional with multiple cameras) joint information from images of people, with resolution options of 15, 18, or 25 joints (for full-body pose estimation).
The framework's \codeword{OpenPoseDemo.exe} application was invoked to perform feature extraction, using the \codeword{pose\_iter\_584000} model.
The configuration used when invoking the application via \codeword{CMD} is shown in Listing \ref{lst:openpose-benchmark-config} below.

\begin{code}[minted language=bash]
OpenPoseDemo.exe ^
    --image_dir input_images ^
    --write_json output ^
    --display 0 ^
    --render_pose 0 ^
    --scale_number 1
\end{code}
\captionof{listing}{OpenPose Configuration for Benchmarking}
\label{lst:openpose-benchmark-config} \vspace{10pt}

\paragraph{Appearance-based \ac{gr}} A self-authored Python script, based on an example from \cite{songDevelopfengGaitRecognition2024}, was used to perform
(1) Silhouette extraction for each frame and
(2) grayscale GEI image creation and output.
The temporal results of these benchmarks are shown in seconds in Table \ref{tab:gait-data-extraction-time-benchmark-results} below.

\begin{scriptsize}
    \begin{xltabular}{\textwidth}{|p{0.25\textwidth}|X|}
        \caption{Gait Data Extraction Time Benchmark Results}
        \label{tab:gait-data-extraction-time-benchmark-results} \\
        \hline
        \textbf{Approach}         & \textbf{Time (s)}     \\ \hline
        \endfirsthead
        \caption*{Gait Data Extraction Time Benchmark Results (Continued)} \\
        \hline
        \textbf{Approach}         & \textbf{Time (s)}     \\ \hline
        \endhead
        \hline
        \endfoot
        \textbf{Model-Based}      & 2986.8 (49.7 minutes) \\ \hline
        \textbf{Appearance-Based} & 3.1                   \\
    \end{xltabular}
\end{scriptsize}

At face-value, these tests show \ac{gei}-based feature extraction to be far more efficient than the model-based alternative.
However, research explored in Chapter \ref{chap:context} suggests model-based approaches are more accurate and reliable if given enough resources (e.g., a more powerful \ac{gpu}).

Because of the desire for this system to run on low-medium powered hardware, the decision was made to proceed with \ac{gei}-based \ac{gr}.
This strays from what was defined in the \ac{ppd}, which proposed a model-based (i.e., skeleton) \ac{gr} approach.
The intended functionality of the overall system however, remains the same.

\subsection{Implementation Technology Stack}
\label{subsec:implementation-technology-stack}

The \acf{poc} solution was primarily built using the Python programming language HTML, and JavaScript, with a number of libraries and frameworks leveraged throughout, enabling:
\ac{cv} and image processing,
math and matrix transformations for statistical analysis,
\ac{ml}-based feature extraction and classification,
and the creation of \ac{ui}s based on designs illustrated in Section \ref{sec:design}.

A detailed list of the technologies used, their purpose, and justifications for their use are shown in Table \ref{tab:tech-stack} below.

\begin{scriptsize}
    \begin{xltabular}{\textwidth}{|p{0.25\textwidth}|X|}
        \caption{Implementation Technology Stack: Purpose and Justification}
        \label{tab:tech-stack}                                                                                                                                                                                                                 \\
        \hline
        \textbf{Technology}          & \textbf{Purpose and Justification}                                                                                                                                                                      \\ \hline
        \endfirsthead
        \caption*{Implementation Technology Stack: Purpose and Justification (Continued)} \\
        \hline
        \textbf{Technology}          & \textbf{Purpose and Justification}                                                                                                                                                                      \\ \hline
        \endhead
        \hline
        \endfoot
        \textbf{Python}              & The primary programming language for backend development. Ideal for rapid development and prototyping.-, as well as math and \ac{ml}                                        \\ \hline
        \textbf{Flask}               & Web framework for building the backend \ac{api} and routing front-end. Flask is lightweight and flexible, allowing for easy routing and handling of requests with minimal overhead.                                                \\ \hline
        \textbf{OpenCV}              & Library for \ac{cv} tasks, including image processing. \ac{cv2} provides robust tools for image manipulation, making it suitable for tasks like silhouette extraction and background subtraction. \\ \hline
        \textbf{TensorFlow}          & Framework for \ac{ml} tasks. TensorFlow supports efficient model training and deployment, particularly for tasks involving neural networks like \acp{cnn}.                                 \\ \hline
        \textbf{EfficientNetB0}      & Pre-trained \ac{cnn} model for feature extraction from images. EfficientNetB0 offers a good balance between performance and computational efficiency, suitable for real-time applications.                   \\ \hline
        \textbf{SQLAlchemy}          & ORM for database interactions. SQLAlchemy simplifies database operations and provides a high-level abstraction for managing database records.                                                           \\ \hline
        \textbf{SQLite}              & Lightweight database for storing identities and gait samples. SQLite is easy to set up and requires no additional server, making it ideal for development and small-scale applications.                 \\ \hline
        \textbf{HTML/CSS/JavaScript} & Technologies for building the frontend \ac{ui}. These are standard web technologies that allow for responsive and interactive web applications.                                                  \\ \hline
        \textbf{Socket.IO}           & Library for real-time, low-overhead communication between back and front ends of both components. Socket.IO enables efficient data transmission, allowing for real-time updates without the overhead of traditional \ac{http} requests.         \\ \hline
        \textbf{Jinja2}              & Templating engine for rendering HTML pages. Jinja2 allows for dynamic content generation in HTML, with variables able to be passed from the back-end to the front-end, making it easier to create interactive web pages.                                                     \\
    \end{xltabular}
\end{scriptsize}

The list of hardware resources used throughout this project was much shorter and doesn't warrant a table, only requiring a laptop computer (details specified in Table \ref{tab:benchmark-system-specs}) for development and testing, an external webcam (Logitech C920) to capture video input, and a tripod (Manfrotto Compact Light) for mounting the webcam at an acceptable height and angle.
